\structure{ОРГАНИЗАЦИОННО-ЭКОНОМИЧЕСКАЯ~ЧАСТЬ}

Экономическое обоснование является важной составляющей инженерного проекта,
позволяющей определить трудоёмкость работ, необходимые ресурсы и ожидаемую
эффективность внедрения разрабатываемого решения.

Целью организационно-экономического раздела является обоснование
целесообразности разработки методики верификации, оценка затрат на её
создание, а также анализ потенциальных областей применения и
перспектив внедрения. Для достижения цели необходимо решить следующие задачи:

\begin{itemize}
    \item определить структуру и этапы разработки методики и
    программного инструмента;
    \item оценить трудоёмкость выполнения работ;
    \item определить состав исполнителей и распределение задач;
    \item рассчитать затраты на реализацию проекта;
    \item проанализировать область применения и практическую
    значимость разработанного решения.
\end{itemize}
